\chapter{Introduction}

\section{Motivation to Learn Neural Representation}

Brain is an amazing organ. From mathematical perspective, a brain implements a very complex (stateful) function which maps from an animal's sensory information to its behavior. For example, the brain can decide what is food and what is not. The brain enables the animal to navigate in space. The brain can even imagine (or simulate) a future that has not happened before. However, we are very far from fully understand how the brain works, or how it has evolved to become what it is. Understanding the brain will lead to an ability for us to build a better machine intelligence, and open up more possibilities in autonomous system.

Recent advancement in artificial intelligence (AI), specifically in artificial neural network (ANN) takes a lot of inspiration from what we currently know about the brain. ANN model accomplishes in solving many tasks which was considered impossible before, \textit{e.g.}, image recognition, video annotation, speech recognition or machine translation. ANN works by constructing a generalized trainable function mapping from an input to an output, with an architecture that imitates a real neural network. For example, in image recognition task, we can train an ANN to map from an image pixels (real number vector) to its label (usually encoded as a vector with only one non-zero element). Because the function is trainable, it can then be trained on tons of data to do these previously unsolvable task.

However, because the connection inside an ANN model is very complicated, we do not understand how the trained model actually solve the problems. In essence, we have created a wonderful system that can solve difficult tasks by mimicking a part of what we know about the brain, but we do not understand both of these system. This lack of knowledge will eventually impede the progress in development of an intelligent system. Therefore, we need more insight to how the brain works. 

The study on the brain starts in the field of Neuroscience. Neuroscientists study how each individual neuron works. They try to map out functions of the brain by areas. They measure what kind of stimulus a sensory neuron response to, or what kind of action is produced when a motor neuron is stimulated. There are a lot of experimental data in Neuroscience. However, it is lacking in a theory that can bind all these experimental observations together. In addition, a lot the theoretical work focuses on modeling the statistics of the experimental data, instead of working in an algorithmic or functional level. Particularly, we need a theory that capture the principle of information processing that happens in the brain regardless of its implementation. For example, \textcite{aj92} explains information processing happens in the retinal as decorrelation of the visual input singal. \textcite{sc06} explains the filtering properties of auditory nerves by applying the efficient coding principle on natural sounds.

It would be nearly impossible to understand the brain at once in an end-to-end fashion, \textit{i.e.} understand the whole process from an animal receiving a sensory information to responding with a certain behavior. Therefore, we try to simplify the understanding the brain by modularizing the problem. We will assume that on the perception side, the brain creates a general internal representation of the world. Then on the behavior side, an action is produced using information from this internal representation. We will assume further that the internal representation of the world on the perception side is created in a hierarchical fashion, and it becomes more abstract as we go up the hierarchy. An evidence supporting our assumption here is, for example, in the primate visual cortex. The first stage of the visual cortex, the primary visual cortex (V1), responses strongly to visual stimuli that look like edges. The secondary visual cortex (V2) response is similar to V1, but also modulated by a more complex property like illusionary contours, or figure ground discrimination. The inferotemporal cortex which is very high up in the visual pathway hierarchy, is associated with a complex object feature like faces. 

The question then arises: what would be a criteria for creating a representation? \textcite{bh61} proposed that the spikes fired by sensory neurons form a neural code that efficiently represents the sensory information. As a consequence, a neural representation should be learnable by optimizing some objective function which measures efficiency of the representation, and we should be able to understand the brain based on an optimization principle.

Past examples of work along this line of idea that inspire our studies presented in this thesis include the followings. \textcite{ob96} proposed sparse coding as a model for simple cells in the primary visual cortex (V1). Sparse coding model objective function says that the neurons in V1 should preserve visual sensory information as much as possible, while at the same time the firing rate of these neurons should be sparse. The model is able to predict the Gabor's filter\footnote{In image processing, Gabor's filter is a linear filter where its kernel is a product of a sinusoidal plane wave and a Gaussian's envelop. It is usually used for edge detection.}-shaped receptive fields\footnote{In neuroscience, the receptive field of a neuron is a region in the sensory space that the neuron responses to.} of neurons in V1 correctly using numerical optimization. \textcite{ky11} proposed an efficient coding model using a population of noisy linear-nonlinear neurons. This model maximizes the mutual information between the input and the output in the present of noise in both the input signal and the output neurons. The result is that in the regime with high output noise, the neuron receptive field has center-surround shape like one of the retinal ganglion cell\footnote{The center-surround receptive fields of retinal ganglion cells are believed to help decorrelated the raw visual signal, so it can pass through the information bottleneck at the optic nerve more efficiently.}. In contrast, in the regime with low output noise, the receptive field has Gabor's filter shape like neurons in V1. In these studies, the properties of neurons emerge naturally from an optimization principle. Furthermore, the studies show that noise-robustness plays an important role in shaping the properties of a neurons. The studies in this thesis follow these inspirations.

We present two studies in this thesis, where we try to learn neural representation of the data using the denoising autoencoder framework. The denoising autoencoder is an unsupervised learning\footnote{In machine learning, unsupervised learning is a family of methods for modeling the data distribution where the data can be unlabeled (as opposed to supervised learning where the data needs to be labeled, and we usually try to train a model to maps each data point to its label).} model for learning the data distribution structure. It assumes that the data points lie on some manifold embedded in a high dimensional space. When we add random noise to a data point, the data point will leave the manifold most of the time. An artificial neural network model is then trained to map the noisy data point back to its original location, and hence denoise the data. After training, the model will contain information about the data manifold, \textit{i.e.} it learns the structure of the data distribution.

In the first study, we extend the sparse coding model by introducing a group structure for the sparseness constraint. While the regular sparse coding model imposes the sparseness constraint on each individual neuron firing rates, group sparse coding model imposes the constraint on groups of neurons. However, previous group sparse coding models presume the group structure a priori instead of learning it from the data, due to an expensive computational cost. We propose that the denoising autoencoder framework can be used to avoid the cost and learn the group structure. We also show that the model with the group structure fits natural image statistics better than the regular sparse coding model.

In the second studies, we try to understand place cells and grid cells. We suspect that these two kinds of neurons play an important role in animal navigation, since their responses are modulated by the animal location. The receptive field of a place cell is simple: it responses to a single location in space. However, the receptive field of a grid cell perplex neuroscientists: it looks like a triangular grid (hence the name grid cell). The question of why grid cells have this kind of receptive field has remained unsolved. We proposed that the place cell and grid cell population form an attractor network which can memorized an animal location, as well as provide an operation for moving the stored location according to the animal velocity. While place cells give us an easy way for reading out the location (because its receptive field is single modal), grid cells act as an analog error-correcting code. Finally, the connection weights of the attractor network can be learned with the denoising autoencoder framework.

\section{Motivation to Learn the Group Structure for Group Sparse Coding}

Sparse coding is a linear generative model -- it tries to represent data as a linear combination of basis functions. For example, in the visual domain, let $I(x, y)$ be an image pixel at the coordinate $x,y$. Sparse coding tries to represent the image as
\begin{equation}
I(x,y) = \sum_k s_k \Phi_k(x,y)
\end{equation}
where $s_k$ is the $k^{\text{th}}$ coefficient of the linear combination and $\Phi_k$ is the $k^{\text{th}}$ basis function. Note that $\Phi_k$ is a function of $(x,y)$ so $\Phi_k$ itself is an image as well. Therefore, combining multiple basis functions linearly is like superimposing multiple images together, with the ratios determined by the coefficients.

Sparse coding objective is not only to find $s_k$'s that represent the data accurately, but also to find a sparse one, \textit{i.e.} the one where the number of non-zero coefficients is as few as possible. The intuition behind this objective is that we view the data as being constructed by ``parts'' where the parts are the basis functions, and we want to use as few ``parts'' as possible.

In computational neuroscience, we can view each coefficient $s_k$ as a neuron firing rate, and the associated basis function $\Phi_k(x,y)$ as the neuron receptive field. The coefficients can be view as a neural code produced by the neuron population, and this neural code is a representation of the input stimulus (since $I$ is determined exactly when all $s_k$'s are known). This can be view as an efficient representation if we define efficiency as sparsity of neuron firing.

If we try to find a set of basis functions that best fits a natural image dataset, we will find the basis functions looks like one in figure \ref{fig:example_basis}. Note that each basis function (each box in the figure) looks like a Gabor's filter, and they vary in location, orientation and scale. The shapes of the basis functions here happen to qualitatively match the receptive fields of neurons in the primary visual cortex (V1), so we believe that the information processing in V1 must be similar to sparse coding.

\begin{figure}[ht]
\begin{center}
\includegraphics[scale=1.0]{basis_sc_3dB.png}
\end{center}
\caption{Example of basis functions when we apply sparse coding on natural image dataset. An image inside each box is each basis function.}
\label{fig:example_basis}
\end{figure}

Although sparse coding model can explain why the receptive fields of neurons in V1 look like Gabor's filters through the efficient coding hypothesis, the model lacks in a temporal dimension, \textit{i.e} the model accepts a still image, $I(x,y)$, as an input while the actual visual system has to deal with a steam of visual information which changes over time, \textit{i.e.} the incoming signal is actually $I(x,y,t)$. This begs a question of how the visual system in the brain represents the temporal dimension of visual information.

One approach is to extend sparse coding such that it accepts a sequence of still images as its input, and perform sparse coding in this new space where time is spatialized \cite{ob03}.
\begin{equation}
I(x,y,t) = \sum_k s_k \Phi_k(x,y,t)
\end{equation}
This yields spatio-temporal receptive fields instead of spatial receptive fields in regular sparse coding. The results in sensible and interesting. However, it would be much more natural to model the dynamics of the visual inputs with a dynamics in the representation, instead of using spatialization.

Several studies deal with the temporal dimension by modeling an image transformation as a Lie's group transport operator on an image data manifold \cite{cb09, sj10}. This method views each still image as a data point living on a manifold embedded in a high dimensional space, where the whole manifold represents the set of all possible images. An image transformation is then an operator which moves one point on the manifold to another point. The studies use a Lie's group operator to model this kind of transformation. It is also possible to consider the transformation being apply on the representation $s_k$'s instead of on the image itself. Let $s$ be a vector whose components are $s_k$'s. The Lie's group transformation can be written as
\begin{equation}
s(t) = \exp\left[ \sum_k c_k \Gamma_k t \right] s(0)
\end{equation}
Here $s(0)$ and $s(t)$ are the representations at time $0$ and $t$, respectively; $\Gamma_k$ is a Lie's group infinitesimal generator matrix and $c_k$ is a linear coefficient. The summation combines many Lie's group generator together. This model has very high capacity in modeling image transformations, but it is difficult and computational expensive the learn the transformation basis $\Gamma_k$'s.

Another study extends sparse coding by using complex numbers instead of real numbers for the coefficient $s_k$ and the basis function $\Phi_k$ \cite{cc08}. The model imposes the sparsity constraint on the coefficient magnitudes while using the coefficient phase rotation to model the image transformations. It is possible to view each complex coefficient and its associated complex basis function as equivalent to 2 real coefficients and 2 real basis functions (separation of the real and the imaginary parts). Hence, the transformation modeled by the phase rotation is equivalent to a rotation between these 2 basis functions. If we expand this complex sparse coding model into a real sparse coding model by doubling the number of coefficients and basis functions, each phase rotation will be equivalent to a Lie's group transformation similar to the following.
\begin{equation}
\Gamma = 
\begin{bmatrix}
    0& -1& 0& 0& \cdots& 0\\
    1&  0& 0& 0& \cdots& 0\\
    0&  0& 0& 0& \cdots& 0\\
    \vdots&&&&&\vdots\\
    0&  0& 0& 0& \cdots& 0
\end{bmatrix}
\end{equation}
where the first 2 indices of this matrix correspond to the 2 coefficients expanded from a single complex coefficient. We can see that the image transformation of the complex sparse coding model is a special case of the Lie's group approach. Because this model is much more constrained it has less capacity, but there is no need to learned the transformation operators.

Seeing the relationship between the complex sparse coding model and the Lie's group transformation model inspire us to find a middle group between these two models. We would like the image transformation to be learned, but with some constraint on the $\Gamma$ matrix structure, so it can be learned more easily. This is our motivation for trying to find a method that can learn the group structure for the sparse coding model as the group structure can be used to constrain the $\Gamma$ matrix.

\section{Motivation to Understand Grid Cell Population as a Neural Network for Analog Error-Correction}

In contrast to the first study where we proposed a method for learning the group structure on the group sparse coding model, and focus its application in the visual domains, the second study focus on a subnetwork in the brain that is responsible for representing an animal location in space and navigation.

Grid cell is a type of neuron found in the entorhinal cortex of many animal species. One well-known property of grid cell is that its response is modulated by the animal location in space, and its receptive field is multi-modal and periodic. In fact, the locations of the modes look like the nodes of a triangular grid (figure \ref{fig:example_grid}). The dependency of grid cell response on the animal location leads to a hypothesis that grid cell must plays an important role in the animal navigation in space. However, the grid shape receptive field perplex many neuroscientists, and the question of why grid cell receptive field has this shape remains unsolved.

\begin{figure}[ht]
\begin{center}
\includegraphics[scale=0.5]{kz/example_grid.png}
\end{center}
\caption{Example illustration of grid cell receptive field. The blue circles represent an animal position in space where an example grid cell response to.}
\label{fig:example_grid}
\end{figure}

In comparison to grid cells, there are other types of neurons whose responses are modulated by the animal location as well. One of these neuron types is called place cell. Place cells are actually discovered before grid cells, and they are found in the hippocampus, which locates next to the entorhinal cortex. The receptive field of a place cell is single modal, as opposed to the multi-modal receptive field of grid cell. Each place cell responses only to a single location in space. Because the place cell activity depends on the animal location, it is believed to play a role in navigation as well. Therefore, the questions arise: why do the brain needs the grid cells which has more complicated receptive fields, while the network of place cells alone should be able to accomplish the navigation task.

We hope to understand the usefulness of grid cell receptive field properties through an efficient coding principle similarly to how the sparse coding model explains the properties of neurons in the primary visual cortex (V1). At first, we hope that the grid cell receptive fields may emerge naturally from an numerical simulation if the model tries to encode location data, and we have the right objective function. We model a population of grid cells as a system which accepts a location as an input, and tries to encode it robustly with a population of noisy neurons. However, this approach does not succeed, and we find the problem more difficult than we have expected, so we take a step back and try to solve a simpler problem.

The problem we consider is comparing 2 location coding schemes: one utilizes grid cells and the other does not. If we can show that the scheme with grid cells works better than the one without, it  provides us an evidence that the grid receptive fields of grid cells help improve location coding performance.

Because grid cells play an important role in animal navigation, we consider the requirements of the navigation task, and the system incorporating grid cells must fulfill these requirements. First, the neuron population responsible for the animal navigation must be able to memorize its current location for a period of long time with high accuracy. Second, it must be able to do path integration, \textit{i.e.} integrates the animal velocity in order to update the animal location. Third, the location memorized must be easily modifiable with information from an external cue, \textit{i.e.} the location memory must be writable. Finally, the system must be robust to noise coming from the neurons, since real neuron firing rates and connections are not reliable.

Due to the requirements above, we believe that the best architecture for implementing the location memory is an attractor network. An attractor network is a neural network with a recurrent connection, and this connection creates fixed points in the activity dynamics. These fixed points can be used to memorize the population activity patterns. Moreover, path integration can be done in an attractor network by perturbing the energy landscape of the fixed points. In our study, we train an attractor network with a denoising autoencoder framework, which is a novel method.

Even though we can train an attractor network to memorize neuron activity patterns, the choice of the encoding from a location (single real number for 1D space or two real number for 2D space) to the activity pattern is still undecided. This is where we will compare the 2 location coding schemes. We call the first scheme ``place-only'' code and the second one ``place-grid'' code. In the first scheme, all neurons are modeled after place cells -- having single modal receptive fields. In the second scheme, some neurons are modeled after place cells, and some neurons are modeled after grid cells -- having periodic receptive fields. Finding that the ``place-grid'' location coding scheme performs better than the ``place-only'' scheming gives us an answer to the question why grid cells exist.

\section{Outline of This Thesis}
Chapter \ref{chap:group_sparse_coding} of this thesis describes the study on learning the group structure for group sparse coding. Chapter \ref{chap:grid_cells} of this thesis describes the study on understanding grid cell network as an analog error-correcting code. Chapter \ref{chap:conclusion} summarizes the contributions of this thesis and ends with a concluding remarks. A short tutorial on neural network and sparse coding for the readers who are not familiar with the subjects can be found in appendix \ref{app:tutorial}.